\chapter{Syntax reference}
\label{ch:reference}

The whole concrete syntax, on a few pages. Whitespace between tokens is ignored, and a \code{\#} starts a comment that runs to the end of the line.

\section{Predicates}

A predicate compares two arithmetic expressions over signal variables.

\begin{table}[h]
\centering\small
\begin{widefig}\begin{tabular}{@{}l l@{}}
\toprule
\textbf{comparisons} & \code{<} \quad \code{<=} \quad \code{>} \quad \code{>=} \quad \code{==} \quad \code{!=}\\
\textbf{arithmetic} & \code{+} \quad \code{-} \quad \code{*} \quad \code{/} \quad \code{\%} (modulo) \quad \code{\textasciicircum} (power) \quad unary \code{-}\\
\textbf{unary functions} & \code{abs} \; \code{sqrt} \; \code{exp} \; \code{ln} \; \code{log} (base 10) \; \code{sin} \; \code{cos} \; \code{tan} \; \code{floor} \; \code{ceil}\\
\textbf{binary functions} & \code{min(a, b)} \quad \code{max(a, b)} \quad (exactly two arguments)\\
\textbf{numbers} & \code{3.14}, \; \code{.5}, \; \code{1e-6}, \; \code{2.5E3}; \; \code{inf} only as an interval upper bound\\
\bottomrule
\end{tabular}\end{widefig}
\end{table}

\noindent A bare \code{=} is rejected with a hint to write \code{==}. There is no \code{pow} function; use the \code{\textasciicircum} operator. \code{min} and \code{max} take exactly two arguments.

\section{Operators}

\begin{table}[h]
\centering\small
\renewcommand{\arraystretch}{1.2}
\begin{widefig}\begin{tabular}{@{}l l l@{}}
\toprule
\textbf{operator} & \textbf{spellings} & \textbf{notes}\\
\midrule
not & \code{not} \quad \code{!} & unary\\
and & \code{and} \quad \code{\&} \quad \code{\&\&} & \\
or & \code{or} \quad \code{|} \quad \code{||} & \\
implies & \code{implies} \quad \code{->} & \\
always & \code{always} \quad \code{globally} \quad \code{G} & bounded window is the worst case\\
eventually & \code{eventually} \quad \code{finally} \quad \code{F} & bounded window is the best case\\
until & \code{until} \quad \code{U} & binary, interval on the keyword\\
next & \code{next} \quad \code{X} & one step, no interval\\
historically & \code{historically} \quad \code{H} & past always\\
once & \code{once} \quad \code{O} & past eventually\\
since & \code{since} \quad \code{S} & binary past, interval on the keyword\\
probability & \code{P>=p (\ldots)} & directions \code{>=} \code{>} \code{<=} \code{<}; $p \in [0,1]$\\
\bottomrule
\end{tabular}\end{widefig}
\end{table}

\noindent There is no previous-step past operator; \code{next} is the only single-step operator. The probabilistic operator has no bracketed form: write \code{P>=0.9 (phi)}, not \code{P[>=0.9](phi)}, and its body must be parenthesized.

\section{Intervals}

A temporal interval is written \code{[lower, upper]} after the operator keyword. For the prefix operators it sits between the keyword and the operand, \code{G[0, 10](x < 5)}; for \code{until} and \code{since} it sits on the keyword between the operands, \code{(x > 0) U[0, 2] (y > 0)}. Omit it and it defaults to \code{[0, inf)}, so \code{F(x > 0)} is \code{F[0, inf](x > 0)}. The lower bound is finite and at least zero; the upper bound is a number or \code{inf}. Bounds may be reals.

\section{Precedence}

Formula operators bind loosest to tightest as: \code{implies} (right-associative), \code{or}, \code{and}, \code{until} (right), \code{since} (right), the temporal prefixes, \code{not} and \code{P}, then predicates and parenthesized groups. Nested temporal operators need parentheses: write \code{G[0, 5](F[0, 2](x > 0))}. Inside a predicate, arithmetic binds \code{+ -} loosest, then \code{* / \%}, then \code{\textasciicircum} (right-associative), then unary minus and atoms.

\section{Worked examples}

Each line parses as written.

\begin{table}[h]
\centering\small
\renewcommand{\arraystretch}{1.3}
\begin{widefig}\begin{tabular}{@{}l l@{}}
\toprule
\textbf{formula} & \textbf{meaning}\\
\midrule
\code{x >= 5} & \code{x} is at least 5 now\\
\code{abs(x) < 5} & the magnitude of \code{x} is under 5\\
\code{x - y > 0} & \code{x} currently exceeds \code{y}\\
\code{G[0, 10](x < 5)} & \code{x} stays below 5 across $[0,10]$\\
\code{F[0, 50](x > 10)} & \code{x} rises above 10 somewhere in $[0,50]$\\
\code{(x > 0) U[0, 2] (y > 0)} & \code{x} holds until \code{y} turns positive, within 2\\
\code{(p > 0) implies F[0, 20](q > 0)} & whenever \code{p}, then \code{q} within 20\\
\code{H[0, 2]((x > 0) implies (y > 0))} & over the last 2, \code{x} positive meant \code{y} positive\\
\code{O[0, 1](x > 0)} & \code{x} was positive within the last 1\\
\code{X(x > 0)} & \code{x} is positive at the next step\\
\code{P>=0.95(G[0, 10](x > 5))} & with prob.\ $\ge 0.95$, \code{x} stays above 5 on $[0,10]$\\
\bottomrule
\end{tabular}\end{widefig}
\end{table}